\subsection{受限和}

设 \((\mathbf{G}_i)_{i\in I}\) 是带算子群的一个族，且对于 \(i\in I\) ，设 \(\mathbf{H}_i\) 是 \(\mathbf{G}_i\) 的子集，使得 \(\prod_{i\in I}\mathbf{G}_i\) 由 \((x_i)_{i\in I}\) 的集合 \(i\in I\) 使得 \(x_i\notin H_i\) 是有限的是一个稳定子群 \(\prod_{i\in I}\mathbf{G}_i\) 等于 \(\prod_{i\in I}\mathbf{G}_i\) 中正规，若 \(I\) 是有限的。它被称为受限和，即 \(\mathbf{G}_i\) 关于 \(\mathbf{H}_i\) 。当对所有 \(i\) 除有限个外 \(\mathbf{H}_i\) 是 \(\mathbf{G}_i\) ，受限和是乘积的正规稳定子群。当对所有 \(i\) ，子群 \(\mathbf{H}_i\) 约化为 \(\mathbf{G}_i\) 的单位元时， \(\mathbf{G}_i\) 关于 \(\mathbf{H}_i\) 的直和简称为 \(\mathbf{G}_i\) 的受限和，有时记为 \(\prod_{i\in I}\mathbf{G}_i\) 的分拆集合的双射。对于所有 \(i_0\in I\) ，映射 \(\iota_{i_0}:\mathrm{G}_{i_0}\to \prod_{i\in I}\mathrm{G}_i\) 上的作用，定义为 \(\iota_{i_0}(x) = (x_i)_{i\in I}\) 的x和y，其中 \(x_{i_0} = x\) 和 \(x_i = e_i\) 中正规，若 \(i\neq i_0\) ，是带算子群的单射同态，称为典范嵌入。 \(\mathbf{G}_i\) 与稳定子群 \(\operatorname {Im}(\iota_i)\) 等同。子群 \(\mathbf{G}_i\) 是正规的。对于 \(i\neq j\) ，的元素 \(\mathbf{G}_i\) 和 \(\mathbf{G}_j\) 交换且 \(\mathrm{G}_i\cap \mathrm{G}_j = \{e\}\) 。群 \(\prod_{i\in I}\mathrm{G}_i\) 由集合生成 \(\bigcup_{i\in I}\mathrm{G}_i\) .

命题12. 设 \((\phi_i: \mathbf{G}_i \to \mathbf{K})_{i \in I}\) 是一族带算子的群同态，使得对所有 \(i \in I\) 和 \(j \in I\) 使得 \(i \neq j\) , \(x \in \mathbf{G}_i\) , \(y \in \mathbf{G}_j\) ，元素 \(\phi_i(x)\) 和 \(\phi_j(y)\) 于 \(\mathbf{K}\) 交换；存在且仅存在一个带算子的群同态 \(\Phi\) 于 \(\prod_{i \in I} \mathbf{G}_i\) 到 \(\mathbf{K}\) 的集合，使得 \(\phi_i = \Phi \circ \iota_i\) 为所有 \(i \in I\) 。对于每个元素 \(x = (x_i)_{i \in I}\) 于 \(\prod_{i \in I} \mathbf{G}_i\) , \(\Phi(x) = \prod_{i \in I} \phi_i(x_i)\) .

上定义了一个处处有定义的合成律。若 \(\Phi\) 和 \(\Phi'\) 是问题的解，它们在 \(\bigcup_{i\in I}G_i\) 上重合，从而在 \(\prod_{i\in I}G_i\) 上重合，由此得到 \(\Phi\) 的唯一性。我们现在证明 \(\Phi\) 的存在性：对每个元素 \(x = (x_i)_{i\in I}\) 于 \(\prod_{i\in I}G_i\) ，设 \(\Phi(x) = \prod_{i\in I}\phi_i(x_i)\) ( \(\S 1\) , 第5号）。显然 \(\Phi \circ \iota_i = \phi_i\) 为所有 \(i\) 和 \(\Phi\) 与位似交换；公式 \(\Phi(xy) = \Phi(x)\Phi(y)\) 由 \(\S 1\) ，第5号，公式(9)。

定义14. 设G是一个带算子的群，且 \((\mathbf{H}_i)_{i\in I}\) 是G的一族稳定子群。若G的每个元素 \((\mathbf{H}_i)\) 且从 \(\mathbf{H}_i\) 与以下每个元素可交换 \(\mathbf{H}_j\) 对于 \(j\neq i\) 到G的唯一同态，其在每个 \(\prod_{i\in I}\mathbf{H}_i\) 上的限制为典范嵌入，则该同态为同构，则称G为子群族 \(\mathbf{H}_i\) 的内受限和（或受限和）。

当I有限时，我们滥用语言，也称内直积（或直积，或积）代替内受限和。G的每个稳定子群H，若存在一个稳定子群 \(H'\) 的G，使得G是H与 \(H'\) 的直积，称为G的一个直因子。

命题13。设G是一个带算子的群，且 \((\mathbf{H}_i)_{i\in I}\) G的一族稳定子群，使得G的每个元素 \(\mathbf{H}_i\) 与以下每个元素可交换 \(\mathbf{H}_j\) 对于 \(j\neq i\) 要使 G 成为子群族 \((\mathbf{H}_i)_{i\in I}\) 的受限和，必要且充分条件是 G 的每个元素 \(x\) 都能唯一地表示为 \(\prod_{i\in I}y_i\) 的x和y，其中 \((y_i)_{i\in I}\) 是 G 中元素的一个有限支撑族，且 \(y_i\in \mathbf{H}_i\) 为所有 \(i\) .

显然。

命题14。设G是一个带算子的群，且 \((\mathbf{H}_i)_{i\in I}\) G 的一个有限稳定子群族。要使 G 成为该子群族的限制和 \((\mathbf{H}_i)\) ，必要且充分的条件是每个 \(\mathbf{H}_i\) 是正规的，且G是商群的乘积 \((\mathbf{G} / \mathbf{H}^i)\) 的x和y，其中 \(\mathbf{H}^i\) 是由 \(\mathbf{H}_j\) 对于 \(j\neq i\) .

该条件是显然必要的。反之，假设G是……的乘积。 \(K_{i} = G/H^{i}\) 并将 G 与 \(K_{i}\) 。则 \(H_{i}\) 被等同于 \(K_{i}\) 的一个子群，因此对于 \(i \neq j\) ，中的每个元素 \(H_{i}\) 与以下每个元素可交换 \(H_{j}\) ；另一方面， \(H^{i}\) 被等同于 \(K_{j}\) 对于 \(j \neq i\) ，因此 \(H_{i} = K_{i}\) 对所有 i 成立，且 G 是 \(H_{i}\) .

命题15。设G是一个带算子的群，且 \((\mathbf{H}_i)_{1\leqslant i\leqslant n}\) G 的一个正规稳定子群序列，使得

\[
\left(\mathrm{H} _ {1} \mathrm{H} _ {2} \dots \mathrm{H} _ {i}\right) \cap \mathrm{H} _ {i + 1} = \{e \} \text { for } 1 \leqslant i \leqslant n - 1,
\]

表示集合 \(\mathbf{H}_1\mathbf{H}_2\ldots \mathbf{H}_n\) 是 \(\mathbf{G}\) ，即 \(\mathbf{H}_i\) .

通过对 \(n\) 进行归纳，这立即归结为对 \(n = 2\) 证明该命题。我们首先证明，若 \(x \in \mathrm{H}_1\) 和 \(y \in \mathrm{H}_2\) , \(x\) 和 \(y\) 是可交换的；因为 \(xyx^{-1}y^{-1} = (xyx^{-1})y^{-1} = x(yx^{-1}y^{-1})\) 且因此（由于 \(\mathrm{H}_1\) 和 \(\mathrm{H}_2\) 是正规的） \(xyx^{-1}y^{-1} \in \mathrm{H}_1 \cap \mathrm{H}_2\) ，即 \(xyx^{-1}y^{-1} = e\) ，由假设可得。此外 \(\mathrm{H}_1\mathrm{H}_2\) 是 \(G\) 的子集，且它在 \(G\) 的位似变换下稳定。由此（根据第3节命题1）可得 \(\mathrm{H}_1\mathrm{H}_2\) 是 \(G\) ，并且容易验证该子群是正规的。最后假设 \(xy = x'y'\) ，其中 \(x \in \mathrm{H}_1\) , \(x' \in \mathrm{H}_1\) , \(y \in \mathrm{H}_2\) , \(y' \in \mathrm{H}_2\) ；于是 \(x'^{-1}x = y'y^{-1}\) ，因此 \(x'^{-1} \in \mathrm{H}_1 \cap \mathrm{H}_2 = \{e\}\) , \(x' = x\) 并且类似地 \(y' = y\) ; \(\mathrm{H}_1\mathrm{H}_2\) 因此是 \(\mathrm{H}_1\) 和 \(\mathrm{H}_2\) .

当所考虑的群是交换群时，使用术语“直和”而非“直积”。

